% !Mode:: "TeX:UTF-8"
% !TEX program = xelatex
%% -*- mode: TeX -*-
%% -*- coding: utf-8 -*-
\documentclass[zihao=-4,b5paper,cn,pad,scheme=chinese,mode=fancy,section,%
  bibend=biber,refsection=chapter,citestyle=gb7714-2015,bibstyle=gb7714-2015,%
  fontset=adobe,result=noanswer]{elegantbook}

% =====================================================================
% Template for writing math books (elegantbook + mathbook.sty).
%
% Usage:
% - Edit the "Book metadata" block below.
% - Put bib entries into references.bib (or change \addbibresource{...}).
% - Add chapters under chapters/ and \include them below.
% - Edit TERMINOLOGY.md and run `make terminology` for the bilingual glossary.
% - Compile: make   (requires XeLaTeX + biber + zhmakeindex in PATH)
%
% Example instance: Jaffe--Taubes, *Vortices and Monopoles*.
% =====================================================================

\ExecuteBibliographyOptions{sorting=nyt}
\usepackage{array}
\usepackage{longtable}
\usepackage{booktabs}
\usepackage{mathbook}
\usepackage[toc,nonumberlist,nopostdot]{glossaries-extra}
\usepackage{glossary-longbooktabs}

% Bilingual terminology table (optional): edit TERMINOLOGY.md, then `make terminology`.
% Omit this block, the \input, and the \printnoidxglossary block if unused.
\newglossary[term-glg]{terms}{term-gls}{term-glo}{中英文术语对照表}
\makenoidxglossaries

\newcommand{\terminologydescription}[2]{%
  \textbf{#1}\newline
  {\small\color{gray}{#2}}%
}

\newglossarystyle{terminologytable}{%
  \setglossarystyle{long-booktabs}%
  \renewenvironment{theglossary}
    {\begin{longtable}{p{0.30\textwidth}p{0.62\textwidth}}}
    {\end{longtable}}%
  \renewcommand*{\glossaryheader}{%
    \toprule
    \textbf{英文术语} & \textbf{推荐译名及说明}\\
    \midrule
    \endhead
    \bottomrule
    \endfoot}%
  \renewcommand*{\glsgroupskip}{}%
}

% This file is generated from TERMINOLOGY.md.
% Run `make terminology` after editing the Markdown source.

\newglossaryentry{term-001}{
  type={terms},
  name={gauge theory},
  sort={gauge theory},
  description={\terminologydescription{规范理论}{经典场论语境；不与“规范分析”混用}}
}

\newglossaryentry{term-002}{
  type={terms},
  name={Yang--Mills--Higgs theory},
  sort={yang--mills--higgs theory},
  description={\terminologydescription{Yang--Mills--Higgs 理论}{可简称 YMH；首次出现写全称}}
}

\newglossaryentry{term-003}{
  type={terms},
  name={gauge potential / connection},
  sort={gauge potential / connection},
  description={\terminologydescription{规范势 / 联络}{二者同指 \(A\)；正文可并列引入}}
}

\newglossaryentry{term-004}{
  type={terms},
  name={curvature / field strength},
  sort={curvature / field strength},
  description={\terminologydescription{曲率 / 场强}{\(F=dA+A\wedge A\)；不单称“场”以免与 Higgs 场混淆}}
}

\newglossaryentry{term-005}{
  type={terms},
  name={Higgs field},
  sort={higgs field},
  description={\terminologydescription{Higgs 场}{也说标量场；取值于内部对称空间}}
}

\newglossaryentry{term-006}{
  type={terms},
  name={gauge group},
  sort={gauge group},
  description={\terminologydescription{规范群}{记为 \(G\)，取矩阵李群}}
}

\newglossaryentry{term-007}{
  type={terms},
  name={internal symmetry space},
  sort={internal symmetry space},
  description={\terminologydescription{内部对称空间}{记为 \(\mathbb{L}\)；\(G\) 作用其上}}
}

\newglossaryentry{term-008}{
  type={terms},
  name={covariant derivative},
  sort={covariant derivative},
  description={\terminologydescription{协变导数}{最小耦合 \(\nabla_A=\nabla+\rho(A)\)}}
}

\newglossaryentry{term-009}{
  type={terms},
  name={Bianchi identity},
  sort={bianchi identity},
  description={\terminologydescription{Bianchi 恒等式}{\(D_AF_A=0\)}}
}

\newglossaryentry{term-010}{
  type={terms},
  name={instanton},
  sort={instanton},
  description={\terminologydescription{瞬子}{不用“瞬时子”}}
}

\newglossaryentry{term-011}{
  type={terms},
  name={vortex},
  sort={vortex},
  description={\terminologydescription{涡旋}{对应金兹堡–朗道方程的解}}
}

\newglossaryentry{term-012}{
  type={terms},
  name={monopole},
  sort={monopole},
  description={\terminologydescription{磁单极子}{对应杨–米尔斯–希格斯方程的解}}
}


% MetaPost (optional): omit all three lines if the book has no mpost figures.
\usepackage[checksum,lineno,labelnames,latex]{mpostinl}
\mpostsetup{template=images/\mpostfilename-#1.mps}
% metapost/mpost-tex.tex — LaTeX preamble for MetaPost figures (mpostinl mposttex)
%
% Loaded from main.tex after \usepackage{mpostinl}.
% Enables Chinese btex/etex labels inside figures.

\begin{mposttex}
  \usepackage{amssymb,bm,xcolor,amsmath}
  \usepackage{CJKutf8}
  \AtBeginDocument{\begin{CJK*}{UTF8}{gkai}}
  \AtEndDocument{\end{CJK*}}
\end{mposttex}

% metapost/mpost-def.tex — project-specific MetaPost macros (mpostinl mpostdef)
%
% Loaded from main.tex after metapost/mpost-tex.tex.
% Put shared figure macros, colors, and helpers here.

\begin{mpostdef}
  numeric u;
  u := 10pt;

  color slash;
  slash := .85white;

  def coordtwo(expr O, len, sx, sy) =
    drawarrow ((-0.1len,0)--(0.9len,0)) shifted O;
    drawarrow ((0,-0.1len)--(0,0.9len)) shifted O;
  enddef;

  def hatchfill expr p =
    fill p
  enddef;
\end{mpostdef}


% =====================================================================
% Book metadata (EDIT THIS BLOCK)
% =====================================================================
\title{涡旋与磁单极子}
\subtitle{静态规范理论的结构}
\extrainfo{这是一本基于作者最新成果的、关于涡旋和磁单极子的相当前沿的严格数学处理著作。}
\institute{哈佛大学·马萨诸塞州剑桥市·02138}
\author{阿瑟·贾菲（Arthur Jaffe）；克利福德·陶布斯（Clifford Taubes）（著）；破旧的大卡车（译）}
\date{1980}
\cover{images/cover.png}

\addbibresource[location=local]{references.bib}

\begin{document}
\thispagestyle{empty}
\addcontentsline{toc}{chapter}{封面}
\maketitle

\frontmatter
\tableofcontents
\addcontentsline{toc}{chapter}{目录}

\mainmatter
\chapter*{前言}
规范理论起源于对电磁现象的描述，其中包括超导现象；同时，它们也已成为基本粒子物理的核心内容，作为弱相互作用、强相互作用以及电磁相互作用的量子理论之出发点。近年来，经典规范理论由于其内在结构而受到拓扑学家与几何学家的广泛关注。物理学家与数学家的相关研究取得了若干重要成果，其中包括对全部“瞬子”（instantons）的发现，以及对其在量子理论中所起作用的部分理解。

本书提出了另一种研究视角。解析方法与拓扑方法的结合，能够为经典规范理论提供有价值的认识，而在这些问题上，单纯的几何方法尚未取得成功。我们尤其运用上述方法，分析了涡旋所对应的金兹堡–朗道方程以及单极子所对应的杨–米尔斯–希格斯方程的一些一般性质；对于这些方程，目前尚未得到闭式形式的解。

第一章和第二章介绍规范理论的基本背景。随后，作者讨论了寻找并刻画涡旋方程与单极子方程解的若干具体问题。第三章至第五章中所发展的方法具有较强的一般性，其适用范围显然不限于本书所研究的方程。第六章则汇集了一些分析工具，这些工具在多种不同类型的问题中都可能发挥作用。

作者感谢 D. Groisser 在研究过程中提出的诸多建议，以及他在校对方面给予的帮助；并特别感谢 Renate D\('\)Arcangelo 在紧张条件下完成的高质量排版工作。正是她的卓越付出，使得本书得以顺利出版。

本书旨在为经典规范理论的数学分析提供一个及时而系统的入门性介绍。作者希望，本书能够同时对理论物理学家和数学家有所裨益。
\begin{flushright}
阿瑟·贾菲（Arthur Jaffe）\\
克利福德·陶布斯（Clifford Taubes）\\
马萨诸塞州剑桥市\\
1980 年 8 月
\end{flushright}

\chapter{规范理论}\label{chap:1}
\begin{refsection}
本导论性章节旨在为后续章节的数学发展提供总体背景说明。我们给出全文所研究的基本方程，并简要介绍它们在物理学中的应用背景。

\section{Yang--Mills--Higgs理论}\label{sec:ymh}

在具有 $d$ 个空间维度的\iemph{Yang--Mills--Higgs}理论中，其动力学变量包括一个\iemph{规范势}（\iemph{联络}）
\[
A = A_j(x)\,dx^j
\]
以及一个\iemph{标量（Higgs）场}
\[
\Phi = \Phi(x) .
\]

在一个时间维加 $d$ 个空间维的情形下，变量形式类似。分量 $A_j(x)$ 取值于一个矩阵李群 $G$ 的李代数 $\mathfrak{g}$，该群 $G$ 被称为\iemph{规范群}。场 $\Phi$ 取值于一个向量空间 $\mathbb{L}$（称为\iemph{内部对称空间}），群 $G$ 作为变换群作用于其上。

Yang--Mills势定义了一个\iemph{场}（\iemph{曲率}）
\begin{equation}\label{eq:1.1.1}
F = dA + A\wedge A = \frac12 F_{ij}(x)\,dx^i\wedge dx^j ,
\end{equation}
其分量为
\begin{equation}\label{eq:1.1.2}
F_{ij}(x)
 = \partial_i A_j(x) - \partial_j A_i(x)
   + [A_i(x),A_j(x)] .
\end{equation}
此处及以下均采用求和约定。

Higgs场通过协变导数与联络发生耦合，这称为“最小耦合”。设 $\rho$ 为群 $G$ 在 $\mathbb{L}$ 上的一个线性表示，则 $\rho$ 诱导出李代数 $\mathfrak{g}$ 在 $\mathbb{L}$ 上的表示，\iemph{协变导数}定义为
\begin{subequations}
  \begin{equation}\label{eq:1.1.3a}
    (\nabla_{A})_j(\Phi) = \nabla_j\Phi + \rho(A_j)(\Phi),
  \end{equation}
  并且
  \begin{equation}\label{eq:1.1.3b}
    D_A\Phi = (\nabla_A)_j(\Phi)\,dx^j .
  \end{equation}
\end{subequations}
在\cref{chap:1}中，我们研究定义在 $\mathbb{R}^2$ 上的\uwave{阿贝尔Higgs模型}\index{Higgs模型!阿贝尔|hyperpage}，其规范群为 $G=U(1)$。
\begin{rmk}
文献引用示例：Belavin 等人的 BPST 瞬子 \cite{belavin1975}；Polyakov 的工作 \cite{polyakov1974,polyakov1977}。
\end{rmk}

\begin{thm}[示例定理]\label{thm:example}
设 $G$ 为紧 Lie 群，$A$ 为 $\mathfrak{g}$-值联络。则 $F_A = dA + A \wedge A$ 满足 Bianchi 恒等式 $D_A F_A = 0$。
\end{thm}

\begin{proof}[Bianchi 恒等式]
直接计算协变外微分即可。
\end{proof}

\printbibliography[heading=subbibliography,title=本章参考文献]
\end{refsection}

\chapter{作图测试}

\section{中文标签 \texttt{dotlabel}}

图中中文须用 \texttt{btex\ldots etex}（不能用引号字符串）；\texttt{metapost/mpost-tex.tex} 已加载 \texttt{CJKutf8}。

\begin{mpostfig}[label=labels-cn-demo,show]
  z0 = origin;
  z1 = (2,1) scaled u;
  z2 = (-1.5,1.5) scaled u;
  draw fullcircle scaled 3u withcolor .7 white;
  for i=0,1,2:
    pickup defaultpen;
    drawoptions(withpen pencircle scaled 3pt);
    dotlabel("", z[i]);
  endfor
  drawoptions();
  dotlabel.rt(btex 圆心 etex, z[0]);
\end{mpostfig}
\chapter{代码高亮示例}

本节演示 \texttt{minted} 插入代码。编译需 \texttt{-shell-escape}，且系统已安装 \texttt{latexminted}（见 \texttt{docs/minted.md}）。

\section{Python}

\begin{minted}{python}
def fib(n: int) -> int:
    a, b = 0, 1
    for _ in range(n):
        a, b = b, a + b
    return a

print(fib(10))
\end{minted}

\section{Wolfram Language（Mathematica）}

Pygments 自带的 \texttt{mathematica} 词法器把所有标识符都标为 \texttt{Name}，\texttt{Integrate}、\texttt{Sin} 等系统函数不会单独上色。本模板提供 \texttt{wolfram} 环境：按 Wolfram 惯例将\textbf{大写开头}的符号视为内置函数（需先执行 \texttt{make minted-setup}，见 \texttt{docs/minted.md}）。

\begin{wolfram}
(* 计算单位球面的高斯曲率 *)
gaussianCurvature[r_] := 1/r^2

Integrate[gaussianCurvature[r], {r, 0, 1}]
\end{wolfram}

行内代码可用 \wl{D[Sin[x]^2, x]}。

对比：下列 \texttt{mathematica} 环境使用 Pygments 默认词法器，内置函数与用户函数外观相同。

\begin{minted}{mathematica}
Integrate[Sin[x], {x, 0, Pi}]
\end{minted}

\section{\LaTeX}

\begin{minted}{latex}
\newcommand{\norm}[1]{\left\lVert #1 \right\rVert}
\[
  \norm{A x} \le \norm{A}\,\norm{x}.
\]
\end{minted}

% \appendix
% \include{chapters/appendixA}

\backmatter
\printindex
\cleardoublepage
\glsaddall[types={terms}]
\markboth{中英文术语对照表}{中英文术语对照表}
\printnoidxglossary[
  type=terms,
  style=terminologytable,
  sort=word,
  title={中英文术语对照表},
  toctitle={中英文术语对照表}
]
\end{document}
